\section*{Chapter \ref{ch:b2:phanerozoic-earth} --- The Phanerozoic Earth}

\begin{solution}{exo:b2:phanerozoic-earth:1}
(a) Paris (\ang{5.1}~N) and New York (\ang{3.6}~S); Cape Town and Sydney lay near \ang{60}~S.
(b) Cape Town: $61.9 - 33.9 = \ang{28.0}$, $28.0 \times 111.2 = \qty{3114}{km}$; New York:
$40.7 + 3.6 = \ang{44.3}$, \qty{4926}{km}. (c) $3114/300 = \qty{10.4}{km}$ per million years,
\qty{1.0}{cm} a year; $4926/300 = \qty{16.4}{km}$ per million years, \qty{1.6}{cm} a year.
\end{solution}

\begin{solution}{exo:b2:phanerozoic-earth:2}
(a) 72, 120.1, 50.5 and \qty{135.4}{Myr}; mean $(444 - 66)/4 = \qty{94.5}{Myr}$. (b) The
end-Permian (Siberian, \qty{1.9}{Myr} apart), the end-Triassic (CAMP, \qty{0.4}{Myr}) and the
end-Cretaceous (Deccan, \qty{0.5}{Myr}). (c) No: the intervals range from 50 to
\qty{135}{Myr}; a mean exists for any set of dates.
\end{solution}

\begin{solution}{exo:b2:phanerozoic-earth:3}
(a) $\qty{91}{m}/\qty{0.036}{mm.yr^{-1}} = \num{2.5e6}$ years. (b) $3.7/0.036 = 103$ times. (c)
Far from any shore, in clear water, little mud from rivers reached the floor; the rain of
algal shells was the main sediment.
\end{solution}

\begin{solution}{exo:b2:phanerozoic-earth:4}
(a) From the river sandstone to the chalk the deposits deepen upwards: a transgression. The
beach sand above the chalk records a regression. (b) No: it records a rise of relative sea level,
which subsidence of the place, or a drop in the supply of sediment, could also produce. (c)
Finding the same transgression, at the same time, on other continents, or removing the
subsidence by \omterm{def:b2:basins-compared:backstripping}{backstripping}.
\end{solution}

\begin{solution}{exo:b2:phanerozoic-earth:5}
(a) $233/97 = \qty{2.4}{m}$ per million years, \qty{2.4}{\micro m} a year. (b) $3.7/0.0024 =
1540$: today's rise is some 1500 times faster. (c) The backstripped curve: drilling has shown
that cycle charts get the timing of falls right but overstate their size, by a factor of two or
more; New Jersey has its own bias (dynamic topography), hence the estimate of $100 \pm
\qty{50}{m}$.
\end{solution}

\begin{solution}{exo:b2:phanerozoic-earth:6}
(a) $\num{e5} \times \num{e9} \times 2900 = \qty{2.9e17}{kg}$ of basalt, and $\num{2.9e17}
\times 0.0075 = \qty{2.2e15}{kg}$ of carbon dioxide, $\num{2.2e15}/0.04401 = \qty{4.9e16}{mol}$.
(b) $\num{278e-6} \times (44.01/28.97) \times \num{5.1e18} = \qty{2.15e15}{kg}$. (c) One pulse
released about as much as the whole air held in 1750. (d) $\num{e17}/\num{2.15e15} = 46.5$: nearly fifty.
\end{solution}

\begin{solution}{exo:b2:phanerozoic-earth:7}
(a) $2630 \times \frac{\pi}{6} \times (\num{e4})^3 = \qty{1.38e15}{kg}$. (b) $\tfrac12 \times
\num{1.38e15} \times (\num{11.2e3})^2 = \qty{8.7e22}{J}$; at \qty{20}{km.s^{-1}},
\qty{2.76e23}{J}. (c) A year of \qty{47}{TW} is $\num{47e12} \times \num{3.156e7} =
\qty{1.48e21}{J}$: 59 and 186 years.
\end{solution}

\begin{solution}{exo:b2:phanerozoic-earth:8}
(a) Dark colour, fine lamination, no burrows, much organic matter, pyrite; often a shift of the
carbon isotopes and signs of warming. (b) The Toarcian event with the Karoo--Ferrar, OAE~1a with
Ontong Java. (c) Dates carry uncertainties of up to millions of years, and many events and
provinces exist, so chance coincidences are likely; a causal link needs the province to start
just before the event and a tracer of the volcanism, such as mercury, in the same sediments.
\end{solution}

\begin{solution}{exo:b2:phanerozoic-earth:9}
(a) $60.9 + 5.1 = \ang{66.0}$, $66.0 \times 111.2 = \qty{7339}{km}$ in \qty{200}{Myr}:
\qty{36.7}{km} per million years, \qty{3.7}{cm} a year. (b) More than twice the
\qty{1.6}{cm} a year since \qty{300}{Ma}. (c) Only the change of latitude is known; the change of
longitude, free in a palaeomagnetic frame, and any detour add to the distance.
\end{solution}

\begin{solution}{exo:b2:phanerozoic-earth:10}
(a) The fraction of floor of age $t$ is $p(t) = \frac2T(1 - t/T)$. Then $\int_0^T tp\,\mathrm{d}t =
\frac2T\bigl(\frac{T^2}2 - \frac{T^2}3\bigr) = \frac T3$ and $\int_0^T \sqrt t\,p\,\mathrm{d}t =
\frac2T\bigl(\frac23 T^{3/2} - \frac25 T^{3/2}\bigr) = \frac8{15}\sqrt T$. (b) $T = 3 \times 64.2 =
\qty{192.6}{Myr}$; $2500 + \frac8{15} \times 350 \times 13.878 = \qty{5091}{m}$. (c) $T/2 =
\qty{96.3}{Myr}$: \qty{4332}{m}; the fall is \qty{759}{m}, \qty{2}{\percent} less than
\qty{774}{m}. (d) Even model: $\frac23\sqrt\tau = \frac23\sqrt{2\bar t} = 0.943\sqrt{\bar t}$;
triangle: $\frac8{15}\sqrt{3\bar t} = 0.924\sqrt{\bar t}$. The answer depends on the mean age far
more than on the shape.
\end{solution}

\begin{solution}{exo:b2:phanerozoic-earth:11}
(a) In $\Delta s$, only the factor $1 - f^{-1/2}$ depends on $f$; divide by its value at $f = 2$,
$1 - 2^{-1/2} = 0.2929$. (b) $f = (1 - 0.2929\,\Delta s/432)^{-2}$: 1.071 for \qty{50}{m}, 1.239
for \qty{150}{m}; 7 to \qty{24}{\percent} faster than today. (c) $432 \times (1 -
0.8^{-1/2})/0.2929 = \qty{-174}{m}$: the sea would fall by \qty{174}{m}. (d) The age spread
changes only as old floor is subducted and replaced, which takes a time of the order of $\tau$,
about \qty{100}{Myr}; the sea follows over tens of millions of years.
\end{solution}

\begin{solution}{exo:b2:phanerozoic-earth:12}
(a) $2^{-7.4/2} = 0.077$: \qty{92}{\percent} gone. (b) $t = 2\log_2(1/0.19) = \qty{4.8}{Myr}$.
(c) $4.8/0.060 = 80$ times faster. (d) Within a long stage, background extinction alone removes
most species, so a short pulse adds to a large total and is hard to separate from it; the
incompleteness of the record also smears last appearances backwards.
\end{solution}

\begin{solution}{pb:b2:phanerozoic-earth:1}
\textbf{1.} 2500, 4250, 5304 and \qty{6000}{m}.
\textbf{2.} $\tau = 2 \times 64.2 = \qty{128.4}{Myr}$; $2500 + 350\sqrt{128.4} = \qty{6466}{m}$.
\textbf{3.} $\num{293e6}/(65000 \times 128.4) = \qty{35.1}{km}$ per million years,
\qty{35.1}{mm} a year: a slow rate.
\textbf{4.} $\frac1\tau\int_0^\tau(2500 + 350\sqrt t)\,\mathrm{d}t = 2500 + \frac23 \times
350\sqrt\tau = \qty{5144}{m}$.
\textbf{5.} The square-root law is too deep for floor older than a few tens of millions of years,
which flattens; the real ages are not evenly spread, young floor being the most abundant; and
sediment, \omterm{def:b2:oceanic-islands:seamount}{seamounts} and plateaus raise the floor.
\textbf{6.} $\num{293e6} \times 5.144 = \qty{1.51e9}{km^3}$, \qty{13}{\percent} more than the
whole ocean: the model is too deep.
\textbf{7.} $v = \qty{17.6}{km}$ per million years; $17.6 \times 128.4 = \qty{2260}{km}$ on
each side.
\textbf{8.} $t = 1000/17.6 = \qty{56.8}{Myr}$; $2500 + 350\sqrt{56.8} = \qty{5138}{m}$.
\textbf{9.} On one flank, with $x = vt$: $\int_0^{v\tau}\bigl(d(\tau) - d(x/v)\bigr)\mathrm{d}x
= vc\int_0^\tau(\sqrt\tau - \sqrt t)\,\mathrm{d}t = \frac13 vc\,\tau^{3/2}$; twice that for both.
$\frac23 \times 17.6 \times 0.350 \times 128.4^{3/2} = \frac23 \times 17.6 \times 0.350 \times
1454.9 = \qty{5975}{km^2}$.
\textbf{10.} $65000 \times 5975 = \qty{3.88e8}{km^3}$.
\textbf{11.} $\num{3.88e8}/\num{1.33e9} = 0.29$.
\textbf{12.} Its floor is younger on average, so hotter, less dense and higher; the band of young
floor along the axis is wider.
\textbf{13.} $\tau = \qty{64.2}{Myr}$; $2500 + \frac23 \times 350\sqrt{64.2} = \qty{4370}{m}$.
\textbf{14.} $5144 - 4370 = \qty{774}{m}$.
\textbf{15.} $\num{293e6} \times 0.774 = \qty{2.27e8}{km^3}$.
\textbf{16.} $\num{2.27e8}/\num{362e6} = \qty{0.627}{km}$, \qty{627}{m}.
\textbf{17.} $\Delta s = 627 \times (3300 - 1025)/3300 = \qty{432}{m}$.
\textbf{18.} $f = \bigl(1 - 0.2929 \times 100/432\bigr)^{-2} = 1.15$: \qty{15}{\percent} faster.
\textbf{19.} $0.149 \times 148 = \qty{22.1}{\mega km^2}$.
\textbf{20.} $48.1 + 0.32 \times 7.7 = \qty{50.6}{\percent}$, \qty{74.9}{\mega km^2}.
\textbf{21.} As it rises it covers more area, so the same volume makes a thinner layer
(\cref{prop:b2:phanerozoic-earth:flooding}).
\textbf{22.} Flooded areas at 100, 200, 300 and \qty{400}{m}: 22.1, 41.7, 57.3 and
\qty{71.2}{\mega km^2}. Steps: $(362 + 11.0) \times 0.1 = \qty{37.3e6}{km^3}$, then 39.4, 41.2
and \qty{42.6e6}{km^3}; cumulated 37.3, 76.7, 117.9 and \qty{160.5e6}{km^3}. The volume to place
is \qty{156.4e6}{km^3}, between 300 and \qty{400}{m}: $300 + (156.4 - 117.9)/42.6 \times 100
= \qty{390}{m}$.
\textbf{23.} $38.7 + 0.90 \times 9.4 = \qty{47.2}{\percent}$ of today's land.
\textbf{24.} Doubling the spreading rate would raise the sea by about \textbf{\qty{390}{m}} and
drown nearly half of today's land.
\end{solution}
